\section{Reference fluctuations and compensated quadratic integration}\label{sec:E-reference}

The preceding section reduced each electric sector to an integral against
an explicit positive reference measure. We now prove the integration
bounds used to turn its deterministic energy estimates into bounds on
that integral. The first bound applies at every integer dimension and
will also be useful for the uniform fugacity bound. The sharper estimate
uses the electric coordinate and a polynomial bound on the empirical
primitive. Both estimates refer to normalized expectations; the reference
partition factors must be restored when computing cylinder weights.

\subsection{The normalization and the circle comparison}

Denote each reference partition integral more generally by \(J_n(w)\) with weight \(w(t)\) in \(dt/(2\pi)\); the weight to normalize is \(w_0=2\sin t\). Its mass is
\[
 J_n(w_0)=\prod_{i=0}^{n-1}\frac{2\cdot16^i(i!)^4}{\pi(2i+1)((2i)!)^2}
 =\exp\left(-\tfrac14\log(n+1)+O(1)\right).
\]
Indeed in \(x=\cos t\) use Gram orthogonalization for \(dx/\pi\) with the Legendre polynomials \(P_i=(2^i i!)^{-1}(d/dx)^i(x^2-1)^i\); integration by parts gives leading coefficient \(2^{-i}(2i)!/(i!)^2\) and norm squared \(2/(2i+1)\) for \(dx\). Include the Vandermonde scaling \(2^{2i}\). The asymptotic follows by Stirling (or the exact adjacent-factor ratios plus Wallis/Stirling leading order fixing limit 1 of the factors).

The circle estimate is the finite upper bound obtained from monotonicity of the centered Toeplitz determinants and the strong Szeg\H{o} limit; see \cite[Theorems 1.1 and 2.4]{Simon2005Szego}. We give a Schur-polynomial proof to retain the exact coefficient and finite-size normalization. The power-sum and Schur orthogonality method appears in the trace-moment calculation of Diaconis and Shahshahani \cite{DiaconisShahshahani1994} and its development by Diaconis and Evans \cite[Section 2]{DiaconisEvans2001}. The contraction argument below applies in every degree.

\begin{lemma}[Circle comparison]\label{lem:03-circle}\label{E:circle-comparison}
For a real smooth test \(f(t)=f_0+2\sum_{k\ge1} f_k\cos kt\), with smooth even reflection, put \(w_\pm=2(1\pm\cos t)\). Then
\begin{equation}
 J_n(w_+ e^f)J_n(w_- e^f)e^{-2n f_0}\le \exp\sum_{k\ge1} k |f_k|^2 .
 \label{eq:G-12}
\end{equation}
\end{lemma}
\begin{proof}
Before centering, the product equals the circle Gram determinant of \(e^f\) of size \(2n\): shift its consecutive frequencies to centered half-integers, then split into cosines and sines by evenness. The two evaluation determinants factor \(\sqrt2\cos(t/2)\) respectively \(\sqrt2\sin(t/2)\) per variable times a polynomial Vandermonde with leading coefficients \(2^j\), yielding exactly the two measures (by determinant expansion/integration).

For completeness this circle bound follows already on polynomials \(f\) and then by approximation. In the Vandermonde-squared probability on \(m\) circles (here \(m=2n\)), the centered integral is the squared norm of \(\exp(\sum_{k\ge1} f_k p_k)\) where \(p_k=\sum_{l=1}^m z_l^k\). Symmetric holomorphic polynomials have the orthonormal alternant-quotient basis (strictly increasing nonnegative monomial powers in the determinant numerator, ordinary Vandermonde denominator), by integrating determinants. Its reproducing kernel summed over all degrees, as a formal graded identity, is \(\prod_{i,j}(1-z_i\bar u_j)^{-1}\) by Cauchy–Binet then the elementary Cauchy determinant formula. This kernel equally expands using \(\exp(\sum p_k(z)\bar p_k(u)/k)\). Thus the functions \(\prod_k (p_k/\sqrt k)^{\alpha_k}/\sqrt{\alpha_k!}\) in each degree have synthesis norm at most 1 from coefficient \(\ell^2\) (matrix times its adjoint equals the identity in the orthonormal basis). The indicated exponential series hence has squared norm at most \(\exp(\sum k|f_k|^2)\); on polynomial input this series converges to the function also pointwise/absolutely, or take radial limits. This proves \eqref{eq:G-12}.

\end{proof}

Write \(\mathbb E_0\) for the two-label \emph{normalized}, independent reference law at nonnegative integer dimensions \(\nu_\sigma\le C e^{\gamma L}\). We keep the two coordinate roles explicit. In the physical parameter \(s\in[-A,A]\),
\[
 q_\sigma=Q_\sigma'=X_{\nu,\sigma}-\nu_\sigma\lambda_L(s)\,ds,
 \qquad
 \ell(Q)=(f,q)=-\sum_\sigma\int_{-A}^A f_\sigma'(s)Q_\sigma(s)\,ds.
\]
Here \(\lambda_L(s)\,ds=dt/\pi\), and \(Q_\sigma\) vanishes at both endpoints. The centering measure is the reference equilibrium measure, not the finite-dimensional expected empirical measure. For a test written in the \(s\)-coordinate, define its cosine coefficients by
\[
 f_\sigma(\theta^{-1}(t))
   =f_{\sigma,0}+2\sum_{k\ge1}f_{\sigma,k}\cos(kt).
\]
The squared dual \(h_0\)-norm of \(\ell\) is consequently
\(\|\ell\|_*^2=\sum_{\sigma,k}k|f_{\sigma,k}|^2\).
Regularity of a physical test below refers to \(s\), whereas this norm uses its cosine coefficients in \(t\). The following estimate keeps the reference partition masses visible.
\begin{theorem}[Normalized reference source bound]\label{E:reference-source}
For two independent normalized reference ensembles with dimensions
$0\le\nu_\sigma\le C_0e^{\gamma L}$, $L\ge1$, and every real smooth
reflected test $f_\sigma$, one has
\begin{equation}\label{eq:G-13}
 \mathbb E_0 e^{\ell(Q)}
 \le \prod_\sigma J_{\nu_\sigma}(w_0)^{-1}
       \exp\bigl(\|\ell\|_*^2/2\bigr)
 \le\exp\bigl(C L+\|\ell\|_*^2/2\bigr).
\end{equation}
The constant depends only on $C_0$ and $\gamma$; a zero-dimensional
ensemble consists of the empty configuration and has partition mass one.
\end{theorem}
\begin{proof}
For each label, Cauchy--Schwarz in the unnormalized integral and
$w_0=\sqrt{w_+w_-}$ bound the tilted integral by the geometric mean
of the two tilted $w_\pm$ integrals. Lemma~\ref{lem:03-circle} bounds
that mean after subtracting $\nu_\sigma f_{\sigma,0}$.
Divide by $J_{\nu_\sigma}(w_0)$ and multiply the two estimates.
The exact normalization above gives the last inequality, including
$\nu_\sigma=0$.
\end{proof}
Additive constants in \(f\) do not matter. Use \(\|\cdot\|_{-2,A}\) on primitives for the Sobolev norm with squared weights \((1+(n/A)^2)^{-2}\) in the \(s\)-interval orthonormal sine basis (odd endpoint reflection). These empirical primitives need not be in \(\mathcal X\). Pairings against smooth odd-reflected tests and mixed kernels are nevertheless defined by distributions. On a polynomially bounded \(\|Q\|_{-2,A}\) event, dual Sobolev estimates control these pairings by the physical derivatives of the tests and kernels.

The loss \(CL\) in \eqref{eq:G-13} is suitable for the untruncated count bound, but is of the same order as the exponent we want to determine. The next estimate removes that loss for dimensions comparable to \(m\). Throughout its statement and the sharp quadratic estimate below, polynomial derivative bounds mean the following: for every fixed derivative order \(r\), the relevant \(C^r\)-norm is at most \(C_rL^{a_r}\), with \(C_r,a_r\) fixed before \(L\to\infty\). For a kernel this convention applies to derivatives in both variables. Errors are uniform over families with these fixed bounds.

\begin{theorem}[Uniform source inequality]\label{thm:03-uniform-source}\label{E:sharp-reference-source}
If \(\nu_\sigma\asymp m\), \(\mathcal G=\{\|Q\|_{-2,A}\le L^D\}\) for fixed \(D\), and the tests arise from real smooth even-reflected \(f_\sigma(s)\) with polynomial derivative bounds in the \(s\)-coordinate as specified above, then uniformly,
\begin{equation}
 \mathbb E_0[\mathbf 1_{\mathcal G} e^{\ell(Q)}]
 \le \exp(o(L)+\|\ell\|_*^2/2)
 \label{eq:G-14}
\end{equation}
(where the normalized excess limsup is at most zero). The error is uniform once \(D\), the polynomial derivative bounds, and the constants implicit in \(\nu_\sigma\asymp m\) are fixed.
\end{theorem}
\begin{proof}
We compare the two Jacobi weights \(w_+\) and \(w_-\) after opposite small transports of the particles. Their geometric mean cancels the terms linear in the source and in the empirical fluctuation. The surviving gain will be
\(\gamma\int(a-2a^2)+o(L)\), where \(a\) is the transport primitive. Its maximizing value \(1/4\) gives \(\gamma L/4\) per label, the logarithmic normalization loss up to the small boundary error. We now justify this comparison, including the Jacobian and its uniform error on the event.

\emph{The transport and its smooth error.}
Fix \(0<\delta<1/4\). Choose a smooth primitive \(a\) with
\(0\le a\le1/4\), equal to \(1/4\) on
\(|s|\le(1-2\delta)L\), and zero for \(|s|\ge(1-\delta)L\).
It may be chosen monotone in each transition region, with
\(\|a^{(r)}\|_\infty\le C_{r,\delta}L^{-r}\) for \(r\ge1\).
For each label and \(\epsilon\in\{+1,-1\}\), define
\[
 F_{\sigma,\epsilon}(s)
   =s-\epsilon\frac{a(s)}{\nu_\sigma\lambda_L(s)}.
\]
These are increasing diffeomorphisms, equal to the identity near the
endpoints. For every fixed derivative order,
\[
 \|F_{\sigma,\epsilon}-\mathrm{id}\|_{C^r}
      \le C_{r,\delta}L^{C_r}e^{-c_\delta L}.
\]
Indeed \(\nu_\sigma\lambda_L\ge c e^{\gamma\delta L}\) on the
support of \(a\), and all relative derivatives of \(\lambda_L\)
there are bounded. Write \(m_\sigma=\nu_\sigma\lambda_Lds\).
Taylor expansion of pushforward against a test function gives
\[
 d_{\sigma,\epsilon}:=(F_{\sigma,\epsilon})_*m_\sigma-m_\sigma
       =\epsilon a'\,ds+r_{\sigma,\epsilon}(s)\,ds,
 \qquad
 \|r_{\sigma,\epsilon}\|_{C^r}
       \le C_{r,\delta}L^{C_r}e^{-c_\delta L}.
\]
For example, the first-order term is
\(-[\nu_\sigma\lambda_L(F_{\sigma,\epsilon}-s)]'
=\epsilon a'\); every higher term contains one additional exponentially
small displacement. The remainder is supported in a slightly enlarged
core, has zero integral, and obeys the same bounds after integration.

We prove the comparison for one label, temporarily suppressing
\(\sigma\). Put
\[
 h(s,v)=H_{\rm ref}(\theta(s),\theta(v)),\quad
 K_\epsilon(s,v)=h(F_\epsilon(s),F_\epsilon(v))-h(s,v),\quad
 g(s)=\log\cot(\theta(s)/2).
\]
The ratio of the two cos-differences defining \(K_\epsilon\) extends
smoothly across the diagonal. Its diagonal value is
\[
 K_\epsilon(s,s)
  =-2\log\frac{\sin\theta(F_\epsilon(s))\,
             \theta'(F_\epsilon(s))\,F_\epsilon'(s)}
            {\sin\theta(s)\,\theta'(s)}.
\]
The one-particle base density in \(s\) is
\(w_0(\theta(s))\theta'(s)/(2\pi)
  =\sin\theta(s)\lambda_L(s)\).
Consequently its Jacobian and density ratio cancel exactly the
half-diagonal contribution omitted from the Vandermonde interaction.
For \(X=\sum_i\delta_{s_i}=m+q\), the logarithm of the pulled-back
base-density ratio is therefore exactly
\[
 -\frac12(X,K_\epsilon X).
\]
Since \(hm=0\), its mean--mean and mixed terms satisfy the identities
\[
 (m,K_\epsilon m)=(d_\epsilon,hd_\epsilon),\qquad
 (q,K_\epsilon m)=(q,(hd_\epsilon)\circ F_\epsilon).
\]
It follows that, uniformly on \(\mathcal G\),
\[
 -\tfrac12(X,K_\epsilon X)
       =-\epsilon(q,ha')-\tfrac12h_0(a,a)+o(1).
\]

Here is the precise uniform error estimate. If \(q=Q'\), then
\[
 |(q,b)|\le\|Q\|_{-2,A}\,\|b'\|_{2,A},\qquad
 |(q,Kq)|\le\|Q\|_{-2,A}^2
                \|\partial_s\partial_vK\|_{2,A;2,A},
\]
where the positive Sobolev norms are dual to the negative norm in the
statement, and the last norm is its tensor product. Derivatives through
order three in each variable suffice to bound the last expression,
up to a polynomial factor from the interval length. The divided
difference argument for \(K_\epsilon\), together with the sine bounds
away from the diagonal, bounds all these derivatives by
\(L^Ce^{-c_\delta L}\). The kernel has even reflections at both ends,
so its mixed derivatives have the odd reflections required for this
Sobolev pairing. Convolution of \(h\) with the smooth core-supported
remainder \(r_\epsilon\) has the same bound up to a polynomial factor:
subtract the diagonal logarithm and transfer derivatives to
\(r_\epsilon\), as in Lemma~\ref{lem:03-kernel-regularity}.
The assumptions \(\|Q\|_{-2,A}\le L^D\) and fixed polynomial source
bounds thus make every error in the preceding formula \(o(1)\).
This argument also justifies the expansions that follow.

The change of the source contributes
\[
 (f\circ F_\epsilon-f,X)=\epsilon(f,a')+o(1).
\]
Since \(w_\epsilon/w_0=e^{\epsilon g}\), the extra weight contributes
\[
 \epsilon(g\circ F_\epsilon,X)
       =\epsilon(g,X)+(g,a')+o(1).
\]
Only the core derivatives of \(g\) are involved in this expansion;
although \(g\) diverges at the endpoints, the map is the identity there.
If \(D_\epsilon(X)\) denotes the pulled-back unnormalized density with
weight \(w_\epsilon e^f\), and \(D_0(X)\) the original density with
weight \(w_0e^f\), we have obtained
\[
 \log\frac{D_\epsilon(X)}{D_0(X)}
 =\epsilon\bigl((g,X)+(f,a')-(q,ha')\bigr)
     +(g,a')-\tfrac12h_0(a,a)+o(1).
\]
The terms odd in \(\epsilon\) cancel in the geometric mean. Define
\(G_a=(g,a')-h_0(a,a)/2\). Then on the same event \(\mathcal G\),
\[
 D_0(X)\le e^{-G_a+o(1)}
                 \sqrt{D_+(X)D_-(X)}.
\]

\emph{The gain and the normalization.}
Since \(g'=-\theta'/\sin\theta\),
\(-g'\) is bounded on the core and tends to \(\gamma\) under
translations into either tail whose distance from the endpoints tends
to infinity. Thus
\[
 (g,a')=-\int g'a=\gamma\int a+o_\delta(L).
\]
On the support of \(a'\), two points in the same tail satisfy
\[
 h(s,v)=4\gamma\min(|s|,|v|)
       +O\bigl(1+|\log\min(1,|s-v|)|\bigr),
\]
and opposite-tail values are bounded. To verify this, use
\(\theta(s)\asymp e^{\gamma s}\) in the left tail and
\(\pi-\theta(s)\asymp e^{-\gamma s}\) in the right tail, factor
the difference of the corresponding squared distances, and take its
logarithm. The additive bounded error is uniform on these transition
regions. The logarithmic error has bounded integral against
\(|a'(s)a'(v)|\), since \(\|a'\|_1\le C\) and
\(\|a'\|_\infty\le C_\delta/L\).
For an absolutely continuous \(b\) of compact support on \([0,\infty)\),
\[
 \iint_{[0,\infty)^2}\min(r,t)b'(r)b'(t)\,dr\,dt
       =\int_0^\infty b(r)^2\,dr;
\]
this follows by writing \(\min(r,t)=\int_0^\infty
\mathbf1_{u<r}\mathbf1_{u<t}\,du\). Apply it on each half-line to
obtain
\[
 h_0(a,a)=4\gamma\int a^2+O_\delta(1),\qquad
 G_a=\gamma\int(a-2a^2)+o_\delta(L)
       \ge\frac{\gamma L}{4}-C\delta L-o_\delta(L).
\]

Finally apply the geometric-mean inequality simultaneously to the two
labels, using one common sign \(\epsilon\); this permits the good event
to couple both configurations. Cauchy--Schwarz, removal of that event
from the two positive integrals, and change of variables back to their
original domains give
\[
 \mathbb E_0[\mathbf1_{\mathcal G}e^{\ell(Q)}]
 \le \exp\left(-2G_a+o(1)
             -\sum_\sigma\log J_{\nu_\sigma}(w_0)
             +\tfrac12\|\ell\|_*^2\right)
\]
by \eqref{eq:G-12}; the constant Fourier coefficients cancel because
\(\ell\) is centered. Uniformly for \(\nu_\sigma\asymp m\),
\(\log J_{\nu_\sigma}(w_0)=-\gamma L/4+O(1)\).
Hence the excess over \(\|\ell\|_*^2/2\) is at most
\(C\delta L+o_\delta(L)\). Taking first \(L\to\infty\), then
\(\delta\downarrow0\), proves \eqref{eq:G-14} with the stated uniformity.
\end{proof}
\subsection{Integrating a compensated quadratic cost}

The source inequality bounds linear statistics. The next lemma converts
it into a bound for a coercive quadratic perturbation of the reference
law. It is essential to complete the square before dropping the positive
part of the perturbation: that preserves the correct dependence on the
linear source and the minimizing energy.

\begin{lemma}[Integration of a compensated quadratic cost]\label{lem:04-quadratic}\label{E:quadratic-integration}
Let \(F\) be a bounded real symmetric operator on \(\mathcal X\), with
\(F\ge gI\), \(g>0\), and let \(F-I\) be trace class. Let \(l\in\mathcal X\) be real. Suppose that there are a smooth real vector test \(b(s)\), odd under endpoint reflection, and a smooth real matrix kernel \(\mathscr C(s,v)\) satisfying \(\mathscr C(s,v)=\mathscr C(v,s)^{\mathsf t}\), odd under reflection in either variable, such that
\[
 h_0(l,Q)=\int b(s)^{\mathsf t}Q(s)\,ds,\qquad
 h_0(Q,(F-I)Q)=\iint Q(s)^{\mathsf t}\mathscr C(s,v)Q(v)\,ds\,dv.
\]
Let \(c\in\R\). For the compensated cost
\[
 V(Q)=c+h_0(l,Q)+\tfrac12h_0(Q,(F-I)Q),\qquad
 m_F=c-\tfrac12h_0(l,F^{-1}l),
\]
one has
\begin{equation}
 \log\E_0e^{-V(Q)}
 \le -m_F+CL+\tfrac12\sum_{\lambda<1,\ \lambda\in\operatorname{spec}(F)}
                    -\log\lambda .
 \label{eq:G-18}
\end{equation}
Here and below smooth compensated expressions are evaluated on empirical
primitives by distributional pairing; their individual singular
\(h_0\)-energies are not defined. For the sharp form used below, assume \(\nu_\sigma\asymp m\),
\(g^{-1}+\|F\|+\|l\|_{\mathcal X}\le CL^a\) for fixed \(C,a\), and polynomial physical derivative bounds for \(b\) and \(\mathscr C\) in the convention above. Then for every fixed \(D\), \eqref{eq:G-18} holds with \(CL\) replaced by \(o(L)\) and expectation restricted to \(\|Q\|_{-2,A}\le L^D\). The error is uniform once these bounds and \(D\) are fixed. The constant term \(c\) is unrestricted: its factor \(e^{-c}\) occurs on both sides of the estimate.
\end{lemma}

\begin{proof}
First suppose that all data are generated by finitely many physical sine
tests. Put \(Q_*=-F^{-1}l\), and complete the square before discarding
any positive quadratic terms:
\[
 -V(Q)=-m_F+h_0(Q_*,Q)-\tfrac12\|Q_*\|^2
       -\tfrac12h_0(Q-Q_*,(F-I)(Q-Q_*)).
\]
With \(J=(I-F)_+\), the last term is at most
\(\tfrac12h_0(Q-Q_*,J(Q-Q_*))\). Let \(g_J\) be a real Gaussian
vector of covariance \(J\). Gaussian linearization and
\eqref{eq:G-13} give
\begin{align*}
 \E_0e^{-V(Q)}
 &\le e^{-m_F-\|Q_*\|^2/2}
   \E_{g_J}\!\left[e^{-h_0(g_J,Q_*)}
                      \E_0 e^{h_0(Q_*+g_J,Q)}\right]\\
 &\le e^{-m_F+CL}\E_{g_J}e^{\|g_J\|^2/2}
   =e^{-m_F+CL}\det(I-J)^{-1/2}.
\end{align*}
The cancellation of the mixed and \(Q_*\)-terms in the second line is
exact. The nontrivial eigenvalues of \(I-J\) are precisely those
eigenvalues of \(F\) below one, which proves the formula.

For general smooth data, truncate in the physical \(s\)-sine basis.
The resulting kernels converge with enough derivatives to converge on
every fixed empirical primitive. Their operators converge in trace norm,
and the linear Riesz vectors converge in \(\mathcal X\); these statements
follow by summing the Fourier coefficient tails and using the bounded
embedding \(\mathcal X\hookrightarrow L^2(ds)\). For sufficiently large
cutoffs the gap is at least \(g/2\), so the infima converge. To control
the determinant term, on any finite-dimensional negative-excess subspace
compare a truncated operator to the corresponding compression of \(F\).
For positive operators bounded below by \(g/2\), the integral identity
for \(\log\), or differentiation along their line segment, bounds the
change of log determinant by \(2g^{-1}\) times the trace-norm error.
Interlacing bounds the negative log sum of a compression by that of
\(F\). Fatou's lemma now proves \eqref{eq:G-18}.

The coarse argument also allows the additional cost
\(-t\|Q\|_{-2,A}^2\). Indeed its operator \(B_A\), defined by
\(h_0(Q,B_AQ)=\|Q\|_{-2,A}^2\), has bounded norm and trace
\(O(L)\): factor it through the bounded \(L^2(ds)\) embedding and
sum \((1+(k/A)^2)^{-2}\). Its physical sine truncations also converge
on empirical fields. Choosing \(2t\|B_A\|\le g/2\) preserves a gap.

For the sharpened estimate, take a physical sine cutoff \(K=L^B\).
Smoothness gives errors smaller than any prescribed inverse power of
\(L\), by choosing a sufficiently high fixed derivative order and then
\(B\). Dual Sobolev estimates make the cost error \(o(1)\) on
\(\|Q\|_{-2,A}\le L^D\); trace-norm error divided by the gap and
norm error in the completed square are also \(o(1)\).
The Riesz vectors of the first \(K\) physical sine tests have a Gram
matrix with inverse-polynomial lower bound. To verify this, test the
linear combination of their functionals on the same sine combination
as a primitive. Its \(L^2\)-pairing is the Euclidean squared norm,
whereas the local derivative and boundary estimate for \(h_0\) bounds
its \(\mathcal X\)-norm by a polynomial times that norm. Duality gives
the claimed lower bound. Consequently every polynomially bounded vector
in this finite-dimensional space has a smooth dual test with polynomial
derivative bounds.

Apply \eqref{eq:G-14} in the displayed Gaussian argument when
\(\|g_J\|\le L^{B'}\). Outside this ball use \eqref{eq:G-13}.
After the factor \(e^{\|g_J\|^2/2}\), the normalized tilted Gaussian
has covariance \(J(I-J)^{-1}\), of polynomial norm, in polynomial
dimension. For covariance norm at most \(L^c\) and dimension at most
\(L^c\), its tail outside radius \(L^{B'}\) is bounded, for large
\(B'\), by \(\exp(-L^{2B'-c}/4)\) times an exponential of a fixed
polynomial. Choose \(B'\) to absorb the factor \(e^{CL}\).
This gives the stated \(o(L)\) loss.
\end{proof}

